\chapter{Wstęp}

Liczby losowe stanowią fundament współczesnych systemów bezpieczeństwa informacji.
Wykorzystywane są do generowania kluczy kryptograficznych, wektorów inicjujących,
identyfikatorów sesji oraz wartości jednorazowych, a także w metodach symulacyjnych
Monte Carlo i w algorytmach randomizowanych. W zastosowaniach kryptograficznych jakość
źródła losowości decyduje o bezpieczeństwie całego systemu: nawet poprawnie
zaimplementowany algorytm szyfrowania nie zapewnia ochrony, jeżeli klucz został
wygenerowany w sposób przewidywalny. Historia kryptografii praktycznej dostarcza wielu
przykładów podatności, których źródłem był nie sam algorytm, lecz wadliwy generator
liczb losowych.

Powszechnie stosowane generatory pseudolosowe, oparte na deterministycznych algorytmach,
wytwarzają ciągi o bardzo dobrych właściwościach statystycznych, jednak ich stan
wewnętrzny jest w pełni określony przez wartość początkową. Ciąg taki jest zatem
odtwarzalny, a jego nieprzewidywalność ma charakter wyłącznie obliczeniowy. Tam, gdzie
wymagana jest losowość o charakterze fizycznym, konieczne jest sięgnięcie do zjawisk
niedeterministycznych zachodzących w sprzęcie. Jednym z najczęściej wykorzystywanych
w układach cyfrowych źródeł tego rodzaju jest jitter, czyli losowa fluktuacja momentów
przejść sygnału zegarowego, wynikająca z szumu termicznego i migotania w elementach
półprzewodnikowych.

Układy programowalne FPGA są dla implementacji generatorów sprzętowych platformą
szczególnie atrakcyjną. Umożliwiają realizację źródła entropii wyłącznie przy użyciu
dostępnych zasobów logicznych, bez konieczności stosowania dedykowanych elementów
analogowych, a jednocześnie pozwalają na swobodną modyfikację architektury oraz na pełną
kontrolę nad rozmieszczeniem struktur w obrębie układu. Najczęściej stosowanym
rozwiązaniem tej klasy jest oscylator pierścieniowy, złożony z zamkniętego w pętlę
łańcucha inwerterów o nieparzystej długości, którego przebieg wyjściowy podlega
kumulującemu się jitterowi. Prostota tej struktury jest jednak pozorna: uzyskana jakość
strumienia bitów zależy nie tylko od samego opisu logiki, lecz również od sposobu
odwzorowania jej w zasobach układu, wzajemnego położenia oscylatorów oraz przyjętego
mechanizmu ekstrakcji bitów z próbkowanego przebiegu.

Celem niniejszej pracy jest zaprojektowanie, implementacja oraz przeprowadzenie analizy
statystycznej generatora bitów losowych opartego na oscylatorach pierścieniowych
w strukturze układu FPGA. Praca nie ogranicza się do wykazania działania pojedynczej
implementacji: zasadniczym jej przedmiotem jest zbadanie, w jaki sposób wybrane decyzje
projektowe wpływają na właściwości statystyczne uzyskiwanego strumienia. Analizie poddano
cztery konfiguracje generatora, różniące się liczbą oscylatorów, sposobem ich
rozmieszczenia w strukturze układu, obecnością operacji różnicy symetrycznej na wyjściach
oraz zastosowaniem korekcji rozkładu metodą von Neumanna. Każdą z nich zbadano trzema
niezależnymi metodami: analizą odchylenia rozkładu bitów, badaniem autokorelacji oraz
analizą dwuwymiarowych map rozkładu słów wielobitowych.

Praca składa się z sześciu rozdziałów. Rozdział drugi przedstawia podstawy teoretyczne
analizy statystycznej ciągów losowych: miary entropii, pojęcie ciągu losowego, aparat
testowania hipotez oraz zastosowane w dalszej części metody badawcze wraz z omówieniem
kondycjonowania strumienia i standaryzowanych pakietów testów. Rozdział trzeci omawia
generatory liczb losowych w technice cyfrowej, ich klasyfikację i model funkcjonalny,
a następnie oscylator pierścieniowy jako źródło entropii, ze szczególnym uwzględnieniem
stochastycznego modelu jitteru oraz zjawisk ograniczających entropię. Rozdział czwarty
opisuje wykonaną implementację: wybór platformy sprzętowej, realizację poszczególnych
modułów w języku VHDL oraz interfejs komunikacyjny służący do transmisji danych do
komputera. Rozdział piąty, stanowiący zasadniczą część badawczą, przedstawia narzędzia
programowe opracowane do akwizycji i analizy danych, parametry badanych konfiguracji
oraz wyniki wszystkich przeprowadzonych badań, wraz z ich interpretacją. Rozdział szósty
zawiera podsumowanie pracy oraz wskazanie kierunków dalszego rozwoju.